
\medterm{Q Arm Of A Chromosome} The longer arm of a chromosome, identified by the letter q. A genetic report may use this label to show where a gene change or missing, extra, or rearranged DNA segment is located.

\medterm{Q Fever} An infection caused by the bacterium Coxiella burnetii, usually caught by breathing dust contaminated by infected cattle, sheep, goats, or their birth products. It often causes sudden fever, severe headache, tiredness, and muscle aches; some people develop pneumonia, hepatitis, or long-term heart-valve infection.

\medterm{Q Wave} A downward movement that may begin the electrical pattern of a heartbeat on an electrocardiogram. It can be normal in some recording positions or suggest previous heart-muscle injury, depending on its size and location.

\medterm{Q-T Interval} The time from the start of electrical activation to the end of recovery of the heart’s lower chambers on an electrocardiogram. An unusually long interval can increase the risk of a dangerous abnormal rhythm.



\medterm{QRS Complex} The group of waves on an electrocardiogram showing the electrical signal that makes the heart’s lower chambers contract. An unusually wide or shaped complex may indicate a conduction problem, abnormal rhythm, or electrolyte disturbance.

\medterm{QT Interval} The time on an electrocardiogram representing electrical activation and recovery of the heart's lower chambers. Because heart rate changes this time, clinicians often use a corrected value when assessing rhythm risk.

\medterm{QT Prolongation} An unusually long QT interval on an electrocardiogram. It may result from an inherited condition, a medicine, or abnormal body salts and can increase the risk of a rapid, dangerous rhythm that may cause fainting or cardiac arrest.

\medterm{Quad Screen (Maternal Serum Screening)} A second-trimester blood screening test that estimates the chance of certain chromosome conditions and neural-tube defects in the fetus. It estimates risk rather than making a diagnosis; an abnormal result may require ultrasound or diagnostic testing.

\synonyms
Maternal Serum Screening

\medterm{Quadrant} One of four sections used to describe a location, especially in the abdomen. The four abdominal quadrants help clinicians record where pain, tenderness, a lump, or another finding is located.

\medterm{Quadrantanopia} Loss of vision in one quarter of the normal visual field, often from damage to the brain pathways carrying sight. The missing area may occur on the same side in both eyes and can affect reading, driving, walking, or safety.

\medterm{Quadrantectomy} Surgery to remove one section of an organ, most often part of the breast containing a tumor plus a margin of healthy tissue. The need for lymph-node surgery or further treatment depends on the tumor and test results.

\medterm{Quadriceps} The four muscles on the front of the thigh that extend the knee; the rectus femoris also assists hip flexion. The quadriceps tendon connects them to the patella and contributes to the extensor mechanism.

\medterm{Quadriceps Femoris} The group of four muscles at the front of the thigh. It straightens the knee for standing, walking, climbing, and kicking; one part also helps bend the hip.

\medterm{Quadriceps Femoris Muscle} The four-muscle group at the front of the thigh that straightens the knee. It is essential for standing, walking, climbing, and rising from a chair; injury can cause pain and difficulty extending the leg.

\medterm{Quadriceps Injury} Damage to the quadriceps muscle or tendon caused by sudden force, overuse, or direct trauma. It may cause
front-of-thigh pain, bruising, weakness, or difficulty extending the knee. Severe pain, a palpable defect, or
inability to bear weight requires assessment for a significant tear or tendon rupture.

\medterm{Quadriceps Tendinitis} Overuse-related inflammation or degeneration of the quadriceps tendon, causing pain above the kneecap that worsens with jumping, running, or resisted knee extension.

\medterm{Quadriceps Tendon Rupture} A tear of the tendon connecting the quadriceps muscle to the patella, causing pain, swelling, and inability to actively extend the knee.

\medterm{Quadriceps Tendon Ruptures} A tear of the tendon connecting the quadriceps muscle to the patella, causing weakness
or inability to extend the knee. Complete tears often require surgical repair; treatment
depends on the extent of injury and the person's functional needs.

\medterm{Quadriparesis} Weakness affecting both arms and both legs. It may result from disease or injury to the neck spinal cord, brain, nerves, or muscles and can impair movement, breathing, and bladder or bowel control.

\medterm{Quadriplegia} Paralysis affecting both arms and both legs, usually with weakness of the trunk, because of damage to the neck portion of the spinal cord or another major nervous-system injury. It can impair breathing and bladder or bowel control.

\medterm{Quadruple Screen Test} A second-trimester blood test that combines four pregnancy-related markers to estimate the chance of selected chromosome conditions and neural-tube defects. It does not diagnose a condition; abnormal results need further assessment.

\medterm{Quadruplet Pregnancy} A pregnancy involving four fetuses. It carries increased risks of preterm birth, fetal
growth restriction, hypertensive and metabolic complications, and maternal morbidity and
therefore requires specialized multidisciplinary antenatal care.

\medterm{Quality Of Life} A person’s experience of physical health, emotional well-being, relationships, independence, and ability to do meaningful activities. Illness and treatment may affect these areas differently for each person.

\medterm{Quantitative Bence-Jones Protein Test} A urine test measuring abnormal light-chain proteins made by some plasma-cell cancers. It can help diagnose or monitor disorders such as multiple myeloma, but results are interpreted with blood tests and other findings.

\medterm{Quantitative Computed Tomography} A CT scan that measures bone density while showing bone structure. It can help assess osteoporosis, unusual bone loss, or fracture risk when standard bone-density testing is unsuitable or more detail is needed.

\medterm{Quantitative Immunoglobulins} A blood test measuring the main classes of antibodies. Low levels may suggest an immune deficiency, while high or uneven levels may occur with infection, inflammation, liver disease, or some blood cancers; results need clinical interpretation.

\medterm{Quantitative Nephelometry Test} A laboratory test that estimates the amount of selected proteins by measuring how much light they scatter. It can measure antibodies, complement proteins, and other proteins involved in immune, kidney, or blood disorders.

\medterm{Quantitative PET} Analysis of a PET scan that estimates how much radioactive tracer is taken up by tissue. Standardized measurements can help compare disease activity or treatment response, but results are affected by the tracer, timing, blood sugar, and scanning method.




\medterm{Quarantine} Temporary separation or movement restrictions for people exposed to an infectious disease while monitoring whether they become ill. It differs from isolation, which separates people who are already infected.

\medterm{Quartan Fever} A fever pattern in which attacks recur about every 72 hours, classically seen with the malaria parasite Plasmodium malariae. Fever timing alone cannot diagnose malaria; testing is needed in anyone with possible exposure and fever.













\medterm{Quetiapine Fumarate} A form of quetiapine, a medicine used for schizophrenia, bipolar disorder, and selected depression treatment. It can cause sleepiness, dizziness on standing, weight or blood-sugar changes, and movement problems.

\medterm{Quick Relievers} Fast-acting asthma medicines used when symptoms suddenly worsen. They quickly relax tightened airway muscles and can ease wheezing, coughing, chest tightness, and breathlessness, but they do not replace regular treatment that controls airway inflammation.

\synonyms
Rescue medicines

\medterm{Quickening} The first time a pregnant person feels fetal movement, often as a flutter or tapping sensation during the second trimester. Timing varies with the pregnancy, body shape, and position of the placenta.


\medterm{Quinalphos} A poisonous insecticide that disrupts nerve-signal control. Significant exposure can cause sweating, excess saliva, vomiting, muscle twitching, weakness, seizures, breathing failure, or death and requires emergency poison treatment.

\medterm{Quinapril} Quinapril is an angiotensin-converting-enzyme inhibitor used for hypertension and heart failure; cough, hyperkalaemia, renal-function changes, angioedema, and fetal toxicity are important risks.


\medterm{Quincke's Disease (Angioneurotic Oedema)} Sudden swelling of deeper layers of the skin or mucosa, commonly affecting the lips, face, tongue, throat, or bowel; airway involvement is an emergency.

\synonyms
Angioneurotic Oedema

\medterm{Quincke's Pulse} Visible flushing and paling of the nail bed or skin with each heartbeat. It may occur when the difference between the upper and lower blood-pressure numbers is unusually wide, especially with severe aortic-valve leakage, but it is not diagnostic alone.

\medterm{Quincke's Sign} Alternating redness and paling of the nail bed or skin with each heartbeat, usually caused by a wide pulse pressure. It can be seen in severe aortic-valve leakage but is not diagnostic by itself.

\medterm{Quinidine} A heart-rhythm medicine that changes electrical conduction in the heart. It is used selectively because it can cause dangerous abnormal rhythms, low blood pressure, low platelets, diarrhea, or other serious side effects.

\medterm{Quinine} A medicine derived from cinchona bark that can treat selected malaria infections. It is used less often because safer options are available and it can cause low blood sugar, ringing in the ears, blood-cell problems, or abnormal heart rhythms.


\medterm{Quinolone} A class of antibiotics that stop bacteria from copying and repairing their DNA. Some treat serious infections, but side effects can include tendon injury, nerve symptoms, abnormal heart rhythms, or severe diarrhea, so use depends on the infection and risks.



\medterm{Quinsy} An abscess beside a tonsil, usually after a throat infection. It can cause severe one-sided throat pain, fever, muffled speech, drooling, difficulty opening the mouth, or breathing difficulty and needs urgent treatment.



\medterm{Quinupristin-Dalfopristin} An intravenous antibiotic combination used rarely for selected serious infections caused by resistant gram-positive bacteria. It can cause muscle or joint pain, liver problems, infusion reactions, and medicine interactions.


\medterm{Quit Smoking (Smoking And Quitting Smoking)} Alternate index label for Smoking and Quitting Smoking.


\medterm{Quincke's Edema} An older term for angioedema, sudden swelling in the deeper layers of the skin or mucous membranes. It commonly affects the lips, face, tongue, or throat; throat swelling or breathing difficulty is an emergency.

\medterm{Quotidian Fever} A fever pattern in which fever rises or attacks occur about every 24 hours. It can occur with malaria and other infections, but the timing of fever alone does not identify the cause.
